\section{数一 Fourier 训练：系数、延拓与选点求和}\label{r2:fourier}
对周期为 $2l$ 的分段 $C^1$ 函数，Fourier 系数为
\begin{align*}
 a_0&=\frac1l\int_{-l}^lf(x)\dd x,\\
 a_n&=\frac1l\int_{-l}^lf(x)\cos\frac{n\pi x}{l}\dd x,\qquad
 b_n=\frac1l\int_{-l}^lf(x)\sin\frac{n\pi x}{l}\dd x.
\end{align*}
级数写为 $a_0/2+\sum_{n\ge1}[a_n\cos(n\pi x/l)+b_n\sin(n\pi x/l)]$。周期为 $2l$ 时不能仍机械使用 $\sin nx$、$\cos nx$。Dirichlet 收敛结论取周期延拓后的 $[f(x^-)+f(x^+)]/2$；单个端点被人为赋的函数值不决定和函数。

\subsection{只给 $(0,l)$ 的函数：先选延拓}
奇延拓的正弦系数为 $b_n=(2/l)\int_0^lf(x)\sin(n\pi x/l)\dd x$，没有余弦项与常数项。偶延拓的余弦系数为 $a_n=(2/l)\int_0^lf(x)\cos(n\pi x/l)\dd x$，常数项仍是 $a_0/2$。同一个半区间函数，两种延拓的端点和函数可能不同；选择哪种延拓应由题目所求的级数类型决定。

\begin{example}{正弦展开的端点不能代原函数值}{数一训练}
将 $f(x)=x$（$0<x<l$）作奇延拓。分部积分给
\[
 b_n=\frac2l\int_0^lx\sin\frac{n\pi x}{l}\dd x
 =\frac{2l(-1)^{n+1}}{n\pi}.
\]
因此内部 $x=\sum_{n\ge1}b_n\sin(n\pi x/l)$。在 $x=l$，周期延拓两侧极限分别为 $l,-l$，平均值为零，级数每一项也为零。不能把左极限 $l$ 直接写成端点处的级数和。
\end{example}

\begin{example}{偶函数展开后选点求平方倒数和}{数一训练}
在 $(-\pi,\pi)$ 上取 $f(x)=x^2$，作 $2\pi$ 周期延拓。奇性使 $b_n=0$，且
\[
 a_0=\frac{2\pi^2}{3},\qquad
 a_n=\frac2\pi\int_0^\pi x^2\cos nx\dd x=\frac{4(-1)^n}{n^2}.
\]
最后一步两次分部积分得到：边界的正弦项为零，余下 $2\pi(-1)^n/n^2$。于是
\[
 x^2=\frac{\pi^2}{3}+4\sum_{n=1}^{\infty}\frac{(-1)^n\cos nx}{n^2}.
\]
端点两侧极限均为 $\pi^2$，取 $x=\pi$ 得
\[
 \sum_{n=1}^{\infty}\frac1{n^2}=\frac{\pi^2}{6}.
\]
取 $x=0$ 得 $\sum_{n\ge1}(-1)^{n+1}/n^2=\pi^2/12$。分离偶数项后还有 $\sum_{k\ge0}(2k+1)^{-2}=\pi^2/8$。每次先确认和函数值，再代点，不能只盯三角项会不会消失。
\end{example}

\paragraph{自测与解答}
若 $f(x)=1$（$0<x<\pi$）作奇延拓，则 $b_n=2[1-(-1)^n]/(n\pi)$，即只有奇数项 $4/[(2k+1)\pi]$。在 $x=\pi/2$，函数连续、和函数为一，因此 $1-1/3+1/5-\cdots=\pi/4$；在 $x=0,\pi$，和函数为零。这个条件收敛级数保持自然次序，不随意重排。
